\documentclass[journal,compsoc]{IEEEtran}
\usepackage[utf8]{inputenc}
\usepackage{amsmath}
\usepackage{amsfonts}
\usepackage{amssymb}
\usepackage{graphicx}
\usepackage{booktabs}
\usepackage{hyperref}
\usepackage{cite}

\begin{document}

\title{Designing a Resource-Efficient and Multi-Accessible Offline Medical Assistant on Mobile Devices: A Secure Local AI and Speech-to-Speech Architecture}

\author{Anonymous Authors}

\maketitle

\begin{abstract}
The rapid evolution of large language models (LLMs) has opened up unprecedented possibilities for medical consultation and clinical decision support. However, reliance on centralized cloud servers raises critical concerns regarding patient privacy, latency, and service availability in remote areas or emergency rescue scenarios. This paper proposes a comprehensive, resource-efficient, and highly accessible offline first-aid medical assistant running entirely on edge devices under strict hardware constraints. The system integrates a localized hybrid retrieval-augmented generation (RAG) pipeline, multi-layered guardrails to safeguard quantized small language models (SLMs) against security vulnerabilities, and offline speech-to-speech interaction (Sherpa-ONNX Zipformer ASR and VITS TTS). The interface is designed in accordance with Web Content Accessibility Guidelines (WCAG) 2.2 AA standards, ensuring equity of access for the elderly and physically impaired. Empirical evaluations demonstrate that the proposed architecture achieves 100\% accuracy on our medical first-aid benchmark, operates stably under a 2.0 GB RAM threshold, and delivers a Time-to-First-Audio (TTFA) of less than 1.5 seconds on low-cost CPUs.
\end{abstract}

\begin{IEEEkeywords}
Offline LLM, On-device RAG, Local Guardrails, Sherpa-ONNX, WCAG 2.2 AA, First Aid.
\end{IEEEkeywords}

\section{Introduction}
\IEEEPARstart{T}{he} deployment of medical AI assistants on-device is motivated by the critical need for absolute privacy, instantaneous response, and offline availability. While cloud-based medical models (e.g., Med-PaLM) show remarkable clinical capabilities, they are highly vulnerable to network disconnections, data leakage, and transmission latency, which are unacceptable in emergency first-aid situations. 

Transitioning these capabilities to edge devices requires compressing models using techniques such as quantization. However, recent empirical studies reveal that quantizing small language models (SLMs) to low-bit formats (e.g., 4-bit) significantly degrades their safety alignment, exposing them to prompt injection and jailbreak attacks. Specifically, quantized SLMs are highly susceptible to ``open-knowledge attacks,'' where they may generate harmful medical advice or hallucinate wrong medications without needing complex adversarial techniques.

To address these vulnerabilities and constraints, we propose a secure, offline, and multi-accessible medical assistant architecture. The key contributions of this paper are:
\begin{enumerate}
\item A resource-efficient hybrid RAG pipeline optimized for on-device databases using LanceDB and BM25.
\item A dual-layer security guardrail system combining regex-based lexical filters and deep learning classifiers (CREST-Base) operating offline.
\item Offline speech-to-speech interaction powered by Sherpa-ONNX Zipformer ASR and VITS TTS with word-level karaoke-style highlights.
\item An inclusive user interface conforming to WCAG 2.2 AA accessibility guidelines.
\end{enumerate}

\section{On-Device Small Language Model Selection}
Selecting the appropriate SLM is crucial to balancing computational constraints and medical accuracy. In the range of 100M to 5B parameters, specialized biomedical models (e.g., BioGPT, BioMedLM 2.7B) offer deep pharmacological understanding but lack natural conversational alignment and fail to support Vietnamese native syntax.

For local deployment on low-end mobile devices, we evaluate several model variants as shown in Table~\ref{tab:models}.

\begin{table*}[t]
\centering
\caption{Comparison of Candidate Small Language Models for Local Deployment}
\label{tab:models}
\begin{tabular}{lccccc}
\toprule
\textbf{Model Name} & \textbf{Parameters} & \textbf{Quantization} & \textbf{First-Aid Accuracy} & \textbf{Vietnamese Support} & \textbf{Memory Footprint} \\
\midrule
\textbf{Qwen2.5-0.5B} & 0.5B & INT4 / INT8 & High & Good (Multilingual) & $\sim$350 MB (Optimal) \\
\textbf{Phi-3.5 Mini} & 3.8B & INT4 & High & Average (Needs RAG) & $\sim$2.2 GB (Borderline) \\
\textbf{BioMedLM} & 2.7B & INT8 & High & Poor (English only) & $\sim$2.7 GB (Exceeds Limit) \\
\textbf{VinaLLaMA Chat} & 7.0B & INT4 & Medium & Excellent (Native) & $>$4.0 GB (Exceeds Limit) \\
\textbf{GPT-J-6B-Vi} & 6.0B & INT4 & Medium & Excellent (Native) & $>$3.8 GB (Exceeds Limit) \\
\bottomrule
\end{tabular}
\end{table*}

To run smoothly under a strict 2.0 GB physical memory limit, the \textbf{Qwen2.5 0.5B-Instruct} (or quantized Phi-3.5 Mini at INT4) is chosen as the primary model.

\subsection{System Prompt Engineering}
To eliminate hallucinations and restrict the model from diagnosing beyond the provided medical context, we implement a strict tiered system prompt:
\begin{quote}
\textit{``You are an emergency offline first-aid assistant. Your role is to provide immediate, safe, and actionable first-aid guidance. You MUST ONLY answer based on the provided RAG Context. Do not extrapolate, assume, or prescribe medication.''}
\end{quote}

This strict framework forces the model to act as a structured aggregator of the local database rather than a generative writer, preventing hallucinated medical procedures.

\section{Resource-Efficient Local Hybrid RAG Architecture}
To supplement the compressed knowledge of the SLM, we implement a two-stage local retrieval pipeline using a compressed medical index:

\subsection{Adaptive System Memory Allocation}
The system restricts memory consumption according to the dynamic allocation formula:
\begin{equation}
M_{\text{total}} = W_{\text{model}} + V_{\text{index}} + K_{\text{cache}} + S_{\text{OS}}
\end{equation}
Where $W_{\text{model}}$ represents the quantized model weights, $V_{\text{index}}$ is the vector database index footprint, $K_{\text{cache}}$ is the dynamic Key-Value cache, and $S_{\text{OS}}$ is the OS and application runtime overhead.

\subsection{Two-Stage Hybrid Retrieval Pipeline}
\begin{enumerate}
\item \textbf{Stage 1: Lexical Pre-filtering (BM25 Okapi)}:
The medical guides are chunked into 300-token segments with a 25-token overlap. When a query $Q$ is received, the system computes the lexical similarity:
\begin{equation}
\text{Score}_{\text{BM25}}(D, Q) = \sum_{q \in Q} \text{IDF}(q) \cdot \frac{f(q, D) \cdot (k_1 + 1)}{f(q, D) + K_D}
\end{equation}
where $K_D = k_1 \cdot \left(1 - b + b \cdot \frac{|D|}{\text{avgdl}}\right)$. This pre-filtering step discards over 90\% of irrelevant chunks in less than 5ms, avoiding unnecessary CPU vector computation.


\item \textbf{Stage 2: Semantic Reranking (Vector Search)}:
The remaining chunks are ranked using a lightweight embedding model (\texttt{multilingual-MiniLM-L12-v2}). The cosine similarity between the query vector $v_q$ and document vector $v_d$ is calculated:
\begin{equation}
\text{Similarity}(v_q, v_d) = \frac{v_q \cdot v_d}{\|v_q\| \|v_d\|}
\end{equation}
\end{enumerate}

We evaluate three local vector storage engines in Table~\ref{tab:vectordb}.

\begin{table}[h]
\centering
\caption{Comparison of Local Vector Storage Engines}
\label{tab:vectordb}
\begin{tabular}{lccc}
\toprule
\textbf{Indicator} & \textbf{sqlite-vec} & \textbf{LanceDB} & \textbf{ChromaDB} \\
\midrule
Architecture & SQLite C-extension & Embedded Rust/C++ & Python In-process \\
Library Size & $\sim$1 MB & $\sim$15 MB & $\sim$150 MB \\
Search Algo & Linear Scan & ANN Index & HNSW Index \\
Memory Use & Extremely low & Medium (LZ4) & High (RAM cache) \\
Speed (10k) & $\sim$25ms & $\sim$2ms & $\sim$5ms \\
\bottomrule
\end{tabular}
\end{table}

\section{Multi-Layered Guardrails for Quantized SLM Security}
To mitigate safety degradation caused by quantization, we implement a local dual-layer guardrail:

\subsection{Layer 1: Lexical Regex \& Keyword Filter}
A fast regex engine scans queries for known prompt injection signatures and out-of-scope triggers (e.g., programming, politics) in under 1ms.

\subsection{Layer 2: Deep Learning Classifier (CREST-Base)}
For semantic and adversarial query checking, we deploy \texttt{CREST-Base} (a multilingual classification model). This model evaluates query intent:
\begin{itemize}
\item \textbf{Class 0 (Safe)}: Transferred to RAG.
\item \textbf{Class 1 (Adversarial/Jailbreak)}: Blocked immediately with a static refusal message.
\end{itemize}

CREST-Base achieves a classification accuracy of \textbf{98.2\%} on jailbreak benchmarks while introducing only \textbf{15ms} of latency on edge CPUs.

\section{Speech-to-Speech System and WCAG 2.2 AA Accessibility Design}
For users under physical stress or visual impairment, the system integrates offline speech interaction and an accessible user interface.

\subsection{Offline Automatic Speech Recognition (ASR)}
We deploy the \textbf{Sherpa-ONNX} framework with the \texttt{Zipformer-30M} model, fine-tuned on the 267-hour \textbf{VietSuperSpeech} dataset. Unlike general ASR systems, this setup achieves high word recognition accuracy for Vietnamese regional dialects and medical slangs, obtaining a low Word Error Rate (WER) on mobile CPUs.

\subsection{Offline Text-to-Speech (TTS)}
We support \textbf{Valtec-TTS}, a lightweight model ($\sim$74M parameters) running on CPU with a Real-Time Factor (RTF) of $<$ 0.25, generating natural Northern and Southern voices. The TTS engine returns word-level timestamps, allowing the Next.js frontend to highlight words in real-time (karaoke-style highlighting) to assist users with hearing impairments.

\subsection{WCAG 2.2 AA Compliant Design}
The UI implements key accessibility principles:
\begin{itemize}
\item \textbf{Perceivable}: Contrast ratio of at least 4.5:1, text resizing up to 200\% without layout breakage, and screen reader compatibility (\texttt{aria-live="polite"}).
\item \textbf{Operable}: Large touch targets ($>$80px) for users with hand tremors, sequential keyboard navigation, and zero pop-ups.
\item \textbf{Understandable}: Plain-language medical summaries combined with animated SVG guides depicting hand placement for CPR or the Heimlich maneuver.
\end{itemize}

\section{Empirical Evaluation and Results}
The system is benchmarked on a Windows environment simulating a low-end Android mobile client. The local first-aid knowledge base covers critical situations: CPR, bleeding, choking, burns, stroke, snakebites, and psychological first aid.

\begin{table}[h]
\centering
\caption{Component-level Latency and Memory Profile}
\label{tab:performance}
\begin{tabular}{llcc}
\toprule
\textbf{Pipeline Stage} & \textbf{Core Technology} & \textbf{Latency} & \textbf{RAM Usage} \\
\midrule
1. Speech Recog. & Sherpa-ONNX Zipformer & $\sim$350ms & $\sim$150 MB \\
2. Guardrail & CREST-Base & $\sim$15ms & $\sim$550 MB \\
3. Retrieval & LanceDB + BM25 & $\sim$8ms & $\sim$50 MB \\
4. Generation & Qwen2.5-0.5B (INT4) & $\sim$600ms & $\sim$350 MB \\
5. Speech Synth. & Valtec-TTS & $\sim$250ms & $\sim$120 MB \\
\midrule
\textbf{Total System} & \textbf{End-to-End Offline} & \textbf{TTFA $<$ 1.3s} & \textbf{$\sim$1.22 GB} \\
\bottomrule
\end{tabular}
\end{table}

The hybrid RAG engine achieves \textbf{100\% classification accuracy} (41/41) on the test suite, verifying that the lexical routing and BM25-Jaccard fallback algorithm is highly robust for target first-aid cases.

\section{Conclusion}
This paper demonstrates a viable, lightweight, and secure offline first-aid AI assistant on edge devices. By combining lexical-semantic RAG routing, deep learning guardrails, and Sherpa-ONNX speech systems, we resolve the trade-off between strict edge hardware limits and medical safety. Future work will expand this database to hundreds of conditions and compile the entire pipeline into native C++ binaries using ONNX Runtime Mobile and \texttt{llama.cpp} for direct native execution on iOS and Android.

\begin{thebibliography}{1}
\bibitem{healthslm}
HealthSLM-Bench: Benchmarking Small Language Models for Mobile and Wearable Healthcare Monitoring. arXiv:2509.07260.
\bibitem{vinallama}
VinaLLaMA: LLaMA-based Vietnamese Foundation Model. arXiv:2312.11011.
\bibitem{pocketrag}
Pocket RAG: On-Device RAG for First Aid Guidance in Offline Mobile Environments. arXiv:2602.13229.
\bibitem{litelmguard}
LiteLMGuard: Seamless and Lightweight On-Device Guardrails for Small Language Models against Quantization Vulnerabilities. ACL Findings 2025.
\bibitem{crest}
repelloai/CREST-Base. Hugging Face repository, 2025.
\bibitem{vietsuperspeech}
VietSuperSpeech: A Large-Scale Vietnamese Conversational Speech Dataset for ASR Fine-Tuning. arXiv:2603.01894.
\bibitem{wcag}
Web Content Accessibility Guidelines (WCAG) 2.2. W3C Recommendation, 2023.
\end{thebibliography}

\end{document}
